\section{Direct quantitative separation under exactness}\label{sec:quantitative}

We give a direct proof of strong pure infiniteness from the exact
perturbation principle.  Throughout this section, $C$ is exact and
satisfies $\propP$.  This is $\propF$ with $n=1$, so
Proposition~\ref{prop:perturbation} is available with that multiplicity
and its fixed constant $0<\sigma<1$.  Corollary~\ref{cor:ideal-comparison}
applies to the whole ideal $\Id(h)$ generated by each positive $h$.
For an exact algebra satisfying $\propF$, the scalar conclusion
\eqref{eq:individual-infiniteness} proved in Section~\ref{sec:localization}
supplies $\propP$ before that section invokes the general theorem.

The quantitative feature of this proof is a fixed lower norm bound in
primitive quotients, retained while a mixed product tends to zero.
We first construct positive elements in two prescribed hereditary
subalgebras with this property.  Spectral clipping then supplies bounded
comparison coefficients, so the small product survives Cuntz
implementation.  The resulting estimates verify the strong-infiniteness
inequalities directly.  Thus this proof uses the comparison tools and
the exact perturbation and localization results in
Sections~\ref{sec:foundations}, \ref{sec:exactness}, and~\ref{sec:localization}.

\begin{lemma}[Separation with a fixed quotient norm bound]
\label{lem:separating-positives}
Let $C$ be an exact $C^*$-algebra satisfying $\propP$, and let
$0<\sigma<1$ be the constant in Proposition~\ref{prop:perturbation}
with $n=1$.  Let $a,b\in C_+$ be contractions, and let $0<r<1$.  Set
\[
 m=\frac{r\sigma^2}{1024}.
\]
For every $0<\ell<1/10$ there are positives
$Q_a\in\Her(a)$ and $Q_b\in\Her(b)$ such that
\begin{align}
 \|(Q_a)_P\|&\ge m &&\text{if }\|a_P\|\ge r,\notag\\
 \|(Q_b)_P\|&\ge m &&\text{if }\|b_P\|\ge r,
 \label{eq:separation-quotient-bounds}\\
 \|Q_aQ_b\|&\le20\sqrt\ell.
 \label{eq:separation-product-bound}
\end{align}
Here and below $P$ ranges over $\Prim(C)$.
\end{lemma}

\begin{proof}
Write the polar decomposition in $C^{**}$ as
\[
 ab=vu,\qquad u=|ab|,\qquad u_*=|(ab)^*|.
\]
All three positive elements $a,b,u$ are contractions.  Choose continuous
functions on $[0,1]$ to separate the small singular values of $ab$
from a spectral corner on which an inverse square root is available.
The small values will give small cross products; on the remaining
corner we will preserve exact cancellation while perturbing the two
sides separately.  Use the following functions:
\begin{itemize}
\item $0\le L\le1$, $L=1$ on $[0,5\ell]$, and $L=0$ on
      $[6\ell,1]$; put $N=1-L$;
\item $0\le g\le1$, $g=0$ on $[0,\ell]$, and $g=1$ on
      $[2\ell,1]$;
\item $0\le R\le1$, $R=0$ on $[0,\ell/2]$, and
      $R(s)=\sqrt{\ell/s}$ for $s\ge\ell$.
\end{itemize}
Linear interpolation on the unspecified intervals gives such functions.
A function letter without an argument denotes evaluation at $u$.
$L$ selects the small singular values, $g$ confines the perturbations
away from zero, and $R$ normalizes $u$ on that region.
Functional calculus and the polar decomposition give
\begin{equation}\label{eq:separation-functional-identities}
 L(u_*)v=vL(u),\qquad N(u_*)=vNv^*,\qquad
 gN=N,\qquad (RuR-\ell)g=0.
\end{equation}
The second identity uses $N(0)=0$; the first remains valid with
$L(0)=1$.  Moreover
\[
 vR=ab\,\frac{R(u)}u\in C,\qquad
 vN=ab\,\frac{N(u)}u\in C,
\]
where both quotient functions are extended continuously by zero at
$s=0$, since $R$ and $N$ vanish near zero.
We use the unitization for scalar terms throughout this proof.

\medskip\noindent\textit{The four columns and their products.}
We first specify conditions under which four columns have the product
estimates we need.  Suppose that orthogonal positive contractions
$e,f\in C$ and invertibles $Z,W\in\widetilde C$ satisfy
\begin{equation}\label{eq:separation-inverse-conditions}
 ge=e,\quad gf=f,\quad Z^*W=1,\quad
 Z=1+gzg,\quad W=1+gwg,\quad
 \|Z\|,\|W\|\le2
\end{equation}
for some $z,w\in C$.  For such a choice, define
\begin{equation}\label{eq:separation-four-columns}
 \begin{array}{ll}
 X_0=aL(u_*),&X=\ell^{-1/2}avRZe,\\[2pt]
 Y_0=bL(u),&Y=\ell^{-1/2}bRWf.
 \end{array}
\end{equation}
These elements belong to $C$: the only polar factor is $vR\in C$,
and all unitized factors occur multiplied by elements of $C$.
For any such admissible choice, the proposed outputs are
\begin{equation}\label{eq:separation-final-positives}
 Q_a=(X_0X_0^*+XX^*)^{1/2},\qquad
 Q_b=(Y_0Y_0^*+YY^*)^{1/2}.
\end{equation}
The norm of each positive output dominates the norms of its two
columns in every primitive quotient.  At the end, a row polar
decomposition will bound the product of these outputs by the four
cross products estimated next.

Put $E_R=RuR-\ell$.  Since $E_Rg=0$ and $f=gf$, we have
\[
 E_RWf=E_Rf+E_Rgwgf=0.
\]
In particular, no commutation of $W$ with $g$ is needed.  Using
$v^*ab=u$ and retaining the indicated order of factors gives
\begin{equation}\label{eq:separation-high-orthogonality}
 X^*Y=\ell^{-1}eZ^*RuRWf=eZ^*Wf=ef=0.
\end{equation}
The other three products obey
\begin{equation}\label{eq:separation-mixed-bounds}
 \|X_0^*Y_0\|\le6\ell,\qquad
 \|X^*Y_0\|\le2\sqrt{6\ell},\qquad
 \|X_0^*Y\|\le2\sqrt{6\ell}.
\end{equation}
Indeed,
\begin{align*}
 X_0^*Y_0&=vL(u)uL(u),\\
 X^*Y_0&=\ell^{-1/2}eZ^*RuL(u),\\
 X_0^*Y&=\ell^{-1/2}vL(u)uRWf.
\end{align*}
Here $\|uL(u)\|\le6\ell$.  For $0\le s\le\ell$,
$\ell^{-1/2}R(s)sL(s)\le\sqrt\ell$; for $s\ge\ell$ the same
expression equals $\sqrt{s}L(s)\le\sqrt{6\ell}$.
Together with \eqref{eq:separation-inverse-conditions}, these scalar
estimates prove \eqref{eq:separation-mixed-bounds}.
We now construct $e,f,Z,W$ so that, in each prescribed primitive
quotient, at least one of the two relevant columns also has a fixed
lower norm bound.

\medskip\noindent\textit{Two orthogonal cores.}
Set $h_0=(u-2\ell)_+$.  Since $h_0\oplus h_0\precsim h_0$,
Lemma~\ref{lem:cut-implementation} supplies a row
$[r_1,r_2]\in M_{1,2}(C)$ such that
\begin{equation}\label{eq:separation-exact-cores}
 r_i^*r_j=\begin{cases}\theta(h_0),&i=j,\\0,&i\ne j,\end{cases}
 \qquad r_1r_1^*+r_2r_2^*\in\Her(h_0),
\end{equation}
where $\theta\colon[0,\infty)\to[0,1]$ is continuous, vanishes near
zero, and equals one on $[2\ell,\infty)$.  Define
\[
 e_0=r_1r_1^*,\qquad f_0=r_2r_2^*.
\]
They are orthogonal positive contractions in $\Her(h_0)$; each belongs
to this hereditary algebra because it is dominated by their sum.
In any primitive quotient with $\|u_P\|\ge4\ell$,
$\|(h_0)_P\|\ge2\ell$, and hence
\begin{equation}\label{eq:separation-core-norms}
 \|(e_0)_P\|=\|(f_0)_P\|=\|\theta((h_0)_P)\|=1.
\end{equation}
This includes equality at the threshold $4\ell$ by the chosen plateau.

Put
\[
 k_e=(e_0-1/2)_+,\qquad k_f=(f_0-1/2)_+,
\]
and let $e=\chi(e_0)$ and $f=\chi(f_0)$, where
$\chi\colon[0,1]\to[0,1]$ is continuous, vanishes near zero, and
is one on $[1/2,1]$.
Then
\begin{equation}\label{eq:separation-core-identities}
 ef=0,\qquad ek_e=k_e,\qquad fk_f=k_f,\qquad ge=e,\qquad gf=f.
\end{equation}
For the last two identities, $gh_0=h_0$ implies that multiplication
by $g$ is the identity on $\Her(h_0)$.
By \eqref{eq:separation-core-norms}, both $(k_e)_P$ and $(k_f)_P$
are nonzero whenever $\|u_P\|\ge4\ell$.

\medskip\noindent\textit{The quotient regions requiring a high column.}
Define
\begin{equation}\label{eq:separation-normalization-data}
 \tau=\frac{r\sqrt\ell}{8},\qquad
 C_1=avR,\quad C_2=bR,\qquad
 h_i=(|C_ig|-\tau)_+,\quad d_i=(|C_iN|-2\tau)_+.
\end{equation}
We first check why a lower bound on the region
$\|(d_i)_P\|\ge\tau$ will recover any signal missed by the
corresponding low column.
Suppose $\|a_P\|\ge r$ and $\|(X_0)_P\|<r/2$.
Since $a=aL(u_*)+aN(u_*)$, we have
$\|(aN(u_*))_P\|>r/2$.
The identity $N(u_*)=vNv^*$ gives $\|(avN)_P\|>r/2$.
On $[0,1]$, the scalar inequality $R^2N^2\ge\ell N^2$ holds:
where $N$ is nonzero, $s>5\ell$ and $R(s)^2=\ell/s\ge\ell$.
Consequently
\[
 (C_1N)(C_1N)^*
 =avR^2N^2v^*a
 \ge\ell avN^2v^*a,
\]
and therefore
\[
 \|(C_1N)_P\|>\frac{r\sqrt\ell}{2}=4\tau.
\]
It follows that $\|(d_1)_P\|>2\tau$.
Similarly, if $\|b_P\|\ge r$ and $\|(Y_0)_P\|<r/2$,
then $\|(bN)_P\|>r/2$.  The order inequality
$bR^2N^2b\ge\ell bN^2b$ gives $\|(d_2)_P\|>2\tau$.
Thus every prescribed quotient either has a low column of norm at
least $r/2$ or belongs to the corresponding perturbation region.

\medskip\noindent\textit{Comparison inputs and normalized left factors.}
We now verify the inputs to Proposition~\ref{prop:perturbation}
on these two regions.
If $(d_i)_P\ne0$, then $C_iN=C_igN$ implies
$\|(C_ig)_P\|>2\tau$, so $(h_i)_P\ne0$.  The same assumption
implies $N_P\ne0$, and therefore $\|u_P\|>5\ell$.
Consequently $(k_e)_P,(k_f)_P\ne0$.
The primitive ideal test and Corollary~\ref{cor:ideal-comparison} give
\begin{equation}\label{eq:separation-core-comparisons}
 d_1\precsim h_1,\quad d_1\precsim k_e,\qquad
 d_2\precsim h_2,\quad d_2\precsim k_f.
\end{equation}

Let $D_i=C_ig$, and use continuous functional calculus in the
unitization to set
\begin{equation}\label{eq:separation-left-normalization}
 F_i=\bigl(\max\{\tau,|D_i^*|\}\bigr)^{-1},\qquad
 p_i=F_iD_i.
\end{equation}
Thus $\|F_i\|\le\tau^{-1}$ and $p_i\in C$.
The functional-calculus intertwining identity
\[
 F_iD_i=D_i\bigl(\max\{\tau,|D_i|\}\bigr)^{-1}
\]
follows first for polynomials in $D_iD_i^*$ and $D_i^*D_i$, and then
by uniform approximation.  In particular,
\begin{equation}\label{eq:separation-normalized-core}
 p_i^*p_i=|D_i|^2\bigl(\max\{\tau,|D_i|\}\bigr)^{-2}\le1,
 \qquad p_i^*p_ih_i=h_i.
\end{equation}
The use of $|D_i^*|$ in the left factor $F_i$ is essential to this
identity.

\medskip\noindent\textit{Two compatible inverse perturbations.}
Fix
\[
 \eta=\frac18,\qquad \eta'=\frac{\sigma\eta}{16}
 =\frac{\sigma}{128}.
\]
Apply Proposition~\ref{prop:perturbation}, with $n=1$, to
\[
 (p,q,h,k,c,T,\alpha)
 =(p_1,e,h_1,k_e,d_1,F_1C_1e,\tau)
\]
at scale $\eta$.  Its hypotheses follow from
\eqref{eq:separation-core-identities},
\eqref{eq:separation-core-comparisons}, and
\eqref{eq:separation-normalized-core}.
For the resulting contraction $y\in C$, put
\[
 Z_1=1+\eta gyg.
\]
Since $\|Z_1-1\|\le\eta<1$, the element $Z_1$ is invertible;
define $W_1=(Z_1^*)^{-1}$.  Since $ge=e$,
the conclusion of the Proposition becomes
\begin{equation}\label{eq:separation-first-test}
 \|(F_1C_1Z_1e)_P\|\ge\sigma\eta
 \qquad\text{if }\|(d_1)_P\|\ge\tau.
\end{equation}
Next apply the same Proposition at scale $\eta'$ to
\[
 (p,q,h,k,c,T,\alpha)
 =(p_2,f,h_2,k_f,d_2,F_2C_2W_1f,\tau).
\]
The element $T$ is in $C$; the Proposition imposes no bound on its
norm.  With the resulting contraction $y'\in C$, set
\[
 W=W_1+\eta'gy'g.
\]
Since $gf=f$, the second conclusion is
\begin{equation}\label{eq:separation-second-test}
 \|(F_2C_2Wf)_P\|\ge\sigma\eta'
 \qquad\text{if }\|(d_2)_P\|\ge\tau.
\end{equation}

We record the inverse estimates, including the factorizations needed
to preserve \eqref{eq:separation-first-test}.
For $a_0\in C$ with $\|a_0\|<1$, the Neumann series gives
\begin{equation}\label{eq:separation-factored-series}
 (1+ga_0g)^{-1}=1+g\left(
 \sum_{k=1}^{\infty}(-1)^k a_0(g^2a_0)^{k-1}\right)g.
\end{equation}
The series of middle factors converges in $C$ with norm at most
$\|a_0\|/(1-\|a_0\|)$.  This formula does not require $g$ to be
a projection.  Applying it to $Z_1^*$ gives
\[
 W_1=1+gw_1g,\qquad \|w_1\|\le\frac17.
\]
Thus $W=1+gwg$ for $w=w_1+\eta'y'\in C$, with
\[
 \|w\|\le\frac17+\frac1{128}<\frac14.
\]
Applying \eqref{eq:separation-factored-series} to $W^*=1+gw^*g$ now
proves that $W$ is invertible.  Define $Z=(W^*)^{-1}$.  The same
series gives
\begin{equation}\label{eq:separation-inverse-bounds}
 Z=1+gzg,\qquad \|z\|\le\frac13,
 \qquad \|Z\|\le\frac43,
 \quad\|W\|\le\frac54,
 \quad\|Z_1\|\le\frac98.
\end{equation}
In particular, all conditions in
\eqref{eq:separation-inverse-conditions} hold.

The second perturbation changes the inverse used in the first column.
To keep the first lower bound independent of $\ell$, we must estimate
this change without the potentially large factor $\|F_1\|$.
The inverse difference has the exact order
\[
 Z-Z_1=-\eta' Zg(y')^*gZ_1.
\]
As $Zg=g(1+zg^2)$, its effect on the first normalized column is
bounded by
\begin{align}
 \|F_1C_1(Z-Z_1)e\|
 &=\eta'\|p_1(1+zg^2)(y')^*gZ_1e\|\notag\\
 &\le\eta'\left(1+\frac13\right)\frac98
 =\frac32\eta'.
 \label{eq:separation-no-normalization-loss}
\end{align}
Since $3\eta'/2\le\sigma\eta/2$,
\eqref{eq:separation-first-test} remains bounded below by
$\sigma\eta/2$ with $Z$ in place of $Z_1$.
Using $\|F_i\|\le\tau^{-1}$ in each primitive quotient, and
$\tau\ell^{-1/2}=r/8$, we obtain
\begin{align}
 \|X_P\|&\ge\frac r8\frac{\sigma\eta}{2}
 &&\text{if }\|(d_1)_P\|\ge\tau,\notag\\
 \|Y_P\|&\ge\frac r8\sigma\eta'
 &&\text{if }\|(d_2)_P\|\ge\tau.
 \label{eq:separation-high-lower-bounds}
\end{align}
The cancellation $\tau\ell^{-1/2}=r/8$ is what removes
$\ell$ from both lower bounds.

\medskip\noindent\textit{Assembling the two positive outputs.}

To check the hereditary membership of $Q_a,Q_b$,
\eqref{eq:separation-four-columns} and the bounds on $Z,W$ give
\[
 X_0X_0^*\le a^2,\quad XX^*\le4\ell^{-1}a^2,\qquad
 Y_0Y_0^*\le b^2,\quad YY^*\le4\ell^{-1}b^2.
\]
By Lemma~\ref{lem:comparison-rules}, the first two summands belong
to $\Her(a)$ and the last two to $\Her(b)$.
Thus $Q_a\in\Her(a)$ and $Q_b\in\Her(b)$.
The earlier coverage reduction and
\eqref{eq:separation-high-lower-bounds}, together with domination of
the column norms by the output norms, prove
\eqref{eq:separation-quotient-bounds} with
\[
 m=\frac r8\sigma\eta'=\frac{r\sigma^2}{1024}.
\]
Here $m\le r/2$ and $\eta'\le\eta/2$, so this single lower bound
applies to all cases.  It is independent of $\ell$.

Finally put $U_a=[X_0,X]$ and $U_b=[Y_0,Y]$ in $M_{1,2}(C)$.
Their range squares are $Q_a^2$ and $Q_b^2$.
The rectangular polar decompositions in the bidual have the form
$U_a=Q_av_a$ and $U_b=Q_bv_b$, where
$v_av_a^*$ and $v_bv_b^*$ are the support projections of $Q_a$
and $Q_b$, respectively.  Hence both identities
\[
 U_a^*U_b=v_a^*Q_aQ_bv_b,\qquad
 Q_aQ_b=v_a(U_a^*U_b)v_b^*
\]
hold.  Since partial isometries have norm at most one, they imply
\[
 \|Q_aQ_b\|=\|U_a^*U_b\|
 \le6\ell+4\sqrt{6\ell}\le20\sqrt\ell,
\]
by \eqref{eq:separation-high-orthogonality} and
\eqref{eq:separation-mixed-bounds}.
This proves \eqref{eq:separation-product-bound}.

The construction also covers zero elements and empty norm regions.
If $h_0=0$, the row in \eqref{eq:separation-exact-cores} may be zero.
Empty perturbation targets permit zero perturbations, and the coverage
argument then assigns every prescribed quotient to a low column.
All inverse operations have either the positive floor $\tau$ or the
convergent series \eqref{eq:separation-factored-series}.
\end{proof}

The norms of the preceding outputs may grow as $\ell$ decreases.
The next observation controls the comparison coefficients using only
the fixed spectral floor $m$.  This is enough for their small product
to survive Cuntz implementation.

\begin{lemma}[Comparison above a spectral floor]
\label{lem:bounded-implementation}
Let $C$ be a $C^*$-algebra, let $Q,c\in C_+$ with $\|c\|\le1$,
and let $m>0$.  If
$c\precsim(Q^2-m^2/4)_+$, then for each $0<\delta<1$ there is
$x\in C$ such that
\begin{equation}\label{eq:separation-bounded-implementation}
 \|(Qx)^*(Qx)-c\|<\delta,\qquad
 \|Qx\|\le2,\qquad \|x\|\le\frac4m.
\end{equation}
\end{lemma}

\begin{proof}
Choose a comparison implementer $v\in C$ with
\[
 \|v^*(Q^2-m^2/4)_+v-c\|<\delta.
\]
Such an implementer lies in $C$ by compression to the two scalar
matrix corners.  On $[0,\|Q\|]$ define the continuous function
\[
 j(s)=\begin{cases}s^{-1}(s^2-m^2/4)_+^{1/2},&s>0,\\0,&s=0,
 \end{cases}
 \qquad x=j(Q)v.
\]
Then $(Qx)^*(Qx)=v^*(Q^2-m^2/4)_+v$ and
$\|Qx\|^2<1+\delta<2$, which in particular gives $\|Qx\|\le2$.
Moreover the scalar inequality
\[
 j(s)^2\le\frac4{m^2}s^2j(s)^2
\]
holds because $j$ vanishes on $[0,m/2]$.
Functional calculus followed by multiplication by $v^*$ and $v$ gives
\[
 x^*x\le\frac4{m^2}(Qx)^*(Qx).
\]
Thus $\|x\|\le(2/m)\|Qx\|\le4/m$, independently of the norm
of the original comparison implementer $v$.
\end{proof}

\begin{proposition}[Direct exact proof]\label{prop:direct-strong-infiniteness}
Let $C$ be an exact $C^*$-algebra satisfying $\propP$.  For every
$a,b\in C_+$ and every $\varepsilon>0$, there are $s,t\in C$ such that
\[
 \|s^*a^2s-a^2\|<\varepsilon,\qquad
 \|t^*b^2t-b^2\|<\varepsilon,\qquad
 \|s^*abt\|<\varepsilon.
\]
Consequently $C$ is strongly purely infinite.
\end{proposition}

\begin{proof}
Fix the constant $\sigma$ in Proposition~\ref{prop:perturbation} with
$n=1$.
Let $a,b\in C_+$ be contractions and let $\varepsilon>0$.
Choose $0<r<1$ with $r^2<\varepsilon/3$ and put
$m=r\sigma^2/1024$.  Next choose $0<\ell<1/10$ so small that
\begin{equation}\label{eq:separation-final-ell}
 \frac{320\sqrt\ell}{m^2}<\frac\varepsilon3.
\end{equation}
Take $Q_a,Q_b$ from Lemma~\ref{lem:separating-positives}.
If $(a^2-r^2)_+$ has nonzero image at $P$, then $\|a_P\|>r$,
so $\|(Q_a)_P\|\ge m$.  Thus $(Q_a^2-m^2/4)_+$ has nonzero
image there.  The primitive ideal test and
Corollary~\ref{cor:ideal-comparison} imply
\begin{equation}\label{eq:separation-final-comparisons}
 (a^2-r^2)_+\precsim(Q_a^2-m^2/4)_+,\qquad
 (b^2-r^2)_+\precsim(Q_b^2-m^2/4)_+.
\end{equation}

Finally choose $0<\delta<1$ such that
\begin{equation}\label{eq:separation-final-delta}
 \delta+(4+\delta)\delta<\frac\varepsilon3.
\end{equation}
Apply Lemma~\ref{lem:bounded-implementation} to the comparisons in
\eqref{eq:separation-final-comparisons}.  We obtain $x_a,x_b\in C$
with
\begin{align*}
 \|(Q_ax_a)^*(Q_ax_a)-(a^2-r^2)_+\|&<\delta,\\
 \|(Q_bx_b)^*(Q_bx_b)-(b^2-r^2)_+\|&<\delta,
\end{align*}
and
\[
 \|Q_ax_a\|,\|Q_bx_b\|\le2,\qquad
 \|x_a\|,\|x_b\|\le\frac4m.
\]
Consequently
\begin{equation}\label{eq:separation-implemented-product}
 \|x_a^*Q_aQ_bx_b\|
 \le\frac{16}{m^2}\|Q_aQ_b\|
 \le\frac{320\sqrt\ell}{m^2}.
\end{equation}

Since $Q_a\in\Her(a)=\overline{aCa}$,
$Q_ax_a\in\overline{aC}$; similarly $Q_bx_b\in\overline{bC}$.
Choose $s,t\in C$ such that
\[
 \|as-Q_ax_a\|<\delta,\qquad
 \|bt-Q_bx_b\|<\delta.
\]
For elements $U,V$ of norm at most $2$, replacing them by elements
within $\delta$ changes $U^*U$ or $U^*V$ by at most
$(4+\delta)\delta$, by expanding the difference and using the
triangle inequality.  Also
$\|a^2-(a^2-r^2)_+\|,\|b^2-(b^2-r^2)_+\|\le r^2$.
Therefore
\begin{align*}
 \|s^*a^2s-a^2\|,
 \ \|t^*b^2t-b^2\|
 &\le r^2+\delta+(4+\delta)\delta<\varepsilon,\\
 \|s^*abt\|
 &\le\frac{320\sqrt\ell}{m^2}+(4+\delta)\delta<\varepsilon,
\end{align*}
by \eqref{eq:separation-final-ell} and
\eqref{eq:separation-final-delta}.
The choice order is $r$, then $m$, then $\ell$, then $\delta$;
the lower bound $m$ has no dependence on the later parameters.

For arbitrary positives $a,b$, put $M=\max\{1,\|a\|,\|b\|\}$.
Apply the preceding result to $a/M,b/M$ with tolerance
$\varepsilon/M^2$.  Multiplication of all three inequalities by
$M^2$ gives the required inequalities for $a,b$ with the same
$s,t\in C$.  Zero elements, including the zero algebra, are included.
These are precisely the three inequalities in \eqref{eq:SPI-definition},
so $C$ is strongly purely infinite.
\end{proof}
